% answers_short.tex —— 2010 年·普通高等学校招生考试·广东卷·理科 · 简短答案速查
% 由 tools/make_short_answers.py 自 answers.tex 的 \gksolhead/\gkanswer 提取，
% 并与 2010/2010_数学真题_广东理_答案要点.md 逐题比对后生成，勿手改（改源 MD 后重跑）。

\begin{gknotice}
\textbf{使用说明}：本节只给最终答案便于快速核对，完整推导见解析版。答案为编写者独立推导，\textbf{非官方标准答案}。
\end{gknotice}

\gkpart{选择题}
\begin{center}\begin{tabular}{ccccccccc}
\toprule
题号 & 1 & 2 & 3 & 4 & 5 & 6 & 7 & 8 \\
\midrule
答案 & D & A & B & C & A & D & B & C \\
\bottomrule
\end{tabular}\end{center}

\gkpart{填空题}
\begin{itemize}[leftmargin=2.6em]
  \item[\textbf{9.}] $(2,+\infty)$。
  \item[\textbf{10.}] $2$。
  \item[\textbf{11.}] $1$。
  \item[\textbf{12.}] $(x+2)^2+y^2=2$。
  \item[\textbf{13.}] $\dfrac14$。
  \item[\textbf{14.}] $\dfrac{9a}8$。
  \item[\textbf{15.}] $\left(\sqrt2,\dfrac{3\pi}4\right)$。
\end{itemize}

\gkpart{解答题}
\begin{itemize}[leftmargin=2.6em]
  \item[\textbf{16.}] （1）$\dfrac{2\pi}3$；（2）$f(x)=4\sin\left(3x+\dfrac\pi4\right)$；（3）$\pm\dfrac{\sqrt5}5$。
  \item[\textbf{17.}] （1）$24$ 件；（2）$P(Y=0)=\dfrac2{13}$，$P(Y=1)=\dfrac{32}{65}$，$P(Y=2)=\dfrac{23}{65}$；（3）$\dfrac{144}{625}$。
  \item[\textbf{18.}] （1）证明；（2）$\dfrac{2\sqrt{29}}{29}$。
  \item[\textbf{19.}] 午餐 $4$ 个单位、晚餐 $3$ 个单位，最低费用 $22$ 元。
  \item[\textbf{20.}] （1）$\dfrac{x^2}{2}+y^2=1$（除去 $(\pm\sqrt2,0)$）；（2）$h=\sqrt3$。
  \item[\textbf{21.}] （1）证明；（2）存在。所有点为以 $A$、$B$ 为对角的轴对齐矩形内的一条线段：两坐标同向排列时在直线 $x+y=\frac{x_1+x_2+y_1+y_2}{2}$ 上，反向排列时在直线 $x-y=\frac{x_1+x_2-y_1-y_2}{2}$ 上。
\end{itemize}
